\documentclass[11pt]{article}
\usepackage[margin=25mm]{geometry}
\usepackage{amsmath,amssymb}
\usepackage{graphicx}
\usepackage{booktabs}
\usepackage{tabularx}
\usepackage{placeins}
\usepackage{url}
\setlength{\emergencystretch}{2em}
\newcommand{\pending}[1]{\textit{[Pending: #1]}}
\newcommand{\figslot}[3]{%
\begin{figure}[!htbp]\centering
\fbox{\begin{minipage}[c][36mm][c]{0.92\linewidth}\centering
\textbf{Figure pending}\\[5pt]#1
\end{minipage}}
\caption{#2}\label{#3}\end{figure}}
\title{Cloud-Optical-Thickness-Assisted Estimation of Surface Solar Irradiance from Geostationary Satellite Observations}
\author{}
\date{}
\begin{document}
\begin{center}
{\LARGE\bfseries Cloud-Optical-Thickness-Assisted Estimation of Surface Solar Irradiance from Geostationary Satellite Observations\par}
\end{center}
\vspace{2mm}
\begin{abstract}
Cloud optical thickness (COT) describes cloud attenuation, but its value for irradiance estimation is uncertain when the underlying multispectral imagery is already available. We assess COT-assisted global horizontal irradiance (GHI) estimation at five stations in Hunan and Hubei, with separate regional training. A cloud-product-supervised retriever supplies local COT fields to a compact network receiving AGRI imagery and solar geometry. Across three random initializations, adding COT reduced station-averaged validation RMSE from 123.980 to 120.240 W m$^{-2}$ in Hunan and from 74.333 to 70.993 W m$^{-2}$ in Hubei, reductions of 3.02\% and 4.49\%, respectively. The reduction reached 22.45\% in the Hunan thick-cloud subset defined by the CPP reference product. Daily irradiation mean absolute error also decreased from 0.5491 to 0.5032 kWh m$^{-2}$ over 30 complete Hunan observation days. However, COT did not outperform matched retrieval features or improve a stronger AGRI residual model. In a separate short-term forecasting comparison, adding forecast-derived COT changed GHI RMSE from 77.150 to 78.155 W m$^{-2}$, with a paired confidence interval spanning zero. These validation results identify a benefit of COT in compact satellite irradiance estimators, concentrated in some cloud conditions, but do not establish an advantage over richer image representations or a corresponding forecasting gain.

\end{abstract}
\section{Introduction}
Satellite irradiance estimation connects cloud observations to the radiation received at a ground station. A multispectral model can learn this mapping directly, but its inputs describe spectral responses rather than a single explicit measure of cloud optical thickness. This motivates a practical question: does adding a retrieved COT representation improve GHI estimation when the original imagery and solar geometry are already available?

Multispectral irradiance models use spectral and spatial information to estimate surface radiation. Chen et al. compared selected spectral inputs and attention-based image representations, and examined their errors under different sky conditions \cite{chen2023}. Physical information can also enter the mapping through explicit constraints: Sasaki et al. addressed ambiguity associated with surface reflection and thin clouds using radiative constraints and scene-dependent estimation \cite{sasaki2024}. These approaches motivate testing a more specific question here: whether a retrieved COT field adds useful structure to an irradiance model that already receives the underlying multispectral observations.

Cloud properties also connect current irradiance estimation to short-term forecasting. Papachristopoulou et al. evaluated cloud-related estimation errors separately from cloud-motion-based forecast performance \cite{pap2024}. Recent studies forecast satellite-derived cloud optical depth and use it in irradiance forecasting: Peng et al. combine optical-depth and cloud-mask forecasts with a radiative-transfer-based irradiance estimator \cite{peng2026}; a 2026 study forecasts COT from Meteosat imagery and evaluates its radiative effect through the cloud modification factor \cite{drosoglou2026}. Cloud-FNO evaluates six-hour forecasts of multiple cloud physical properties, including COT, against an independent test set \cite{zhang2026}. These studies establish active links between cloud-property forecasting and solar applications, but differ in satellite, target, spatial support, horizon, and irradiance endpoint. The question here is narrower: whether retrieved COT from the same AGRI observations adds value to matched station-level GHI estimation, and whether that value is retained when a fixed cloud-image extrapolation chain supplies future inputs.

We examine COT as an additional representation learned under cloud-product supervision. Both the imagery-only and COT-augmented models receive the same local AGRI observations and deterministic geometry. The latter also receives a co-registered COT field produced by a separately trained, frozen retriever. COT therefore adds a supervised representation of the available imagery rather than an independent satellite observation. Its usefulness is assessed by the reduction in GHI error under matched data and training conditions.

We ask whether COT improves GHI estimation relative to multispectral imagery, a comparably sized image-feature branch, and features from the retriever itself. Separate regional experiments assess whether the benefit recurs in Hunan and Hubei, while thickness-stratified comparisons examine its dependence on cloud conditions. Finally, a fixed SimVP forecasting model tests whether the benefit persists when future imagery and COT must be predicted rather than observed.

\section{Study Regions and Data}
\subsection{Regional design and observations}
The Hunan validation cohort contains Sili and Zhujia, and the Hubei validation cohort contains Niushou, Wangzhai, and Mizhuang. We report the regional validation results separately and do not interpret them as zero-shot geographic transfer. Regional scores are interpreted within each label definition and are not pooled as if ground observations had identical sampling.

Figure~\ref{fig:regions} locates the five stations and summarizes the first-to-last dates available for AGRI/GHI samples and CPP-supervised retrieval. The current AGRI/GHI record spans 2 April 2024--30 June 2026 in Hunan and 1 August 2025--14 July 2026 in Hubei. CPP-supervised retrieval support ends on 2 March 2026 in both regions and is sparse within these date ranges; the endpoints do not imply continuous daily coverage.

\begin{figure}[!htbp]\centering
\includegraphics[width=0.98\linewidth]{figures/fig1_study_regions_coverage.pdf}
\caption{Study regions and observed date support. (a,b) Station locations in Hunan and Hubei, plotted as local east--west and north--south offsets from each regional station centroid using station coordinates. (c) First and last available dates for the regional AGRI/GHI samples and CPP-supervised COT retrieval. The date lines indicate the span of available records, not continuous daily coverage; CPP is a sparse retrieval reference.}
\label{fig:regions}\end{figure}

\begin{table}[!htbp]\centering\small
\caption{Regional observations and evaluation support. Counts refer to the completed synchronous GHI cohorts within each region.}\label{tab:data}
\begin{tabularx}{\linewidth}{lXr}\toprule
Region & Stations & Validation samples\\\midrule
Hunan & Sili; Zhujia & 4,479\\
Hubei & Niushou; Wangzhai; Mizhuang & 2,456\\\bottomrule
\end{tabularx}
\par\smallskip\footnotesize Hunan: five-minute interval-end observations aggregated to 15 minutes; 2,128/2,351 validation samples at Sili/Zhujia and 52 shared target dates. Hubei: 856/784/816 validation samples at Niushou/Wangzhai/Mizhuang. Its labels use equally weighted available one-minute bins over the right-closed 15-minute interval; missing bins are masked. Hubei irradiance values are interpreted in W m$^{-2}$; source-unit documentation and independent calibration information were unavailable. Regional date partitions and train-only normalization are independent. Because ground-label sampling differs, regional comparisons are within-region.
\end{table}
The satellite input contains thirteen AGRI channels, C01--C06 and C09--C15. Physical channel values are taken from the region-specific geolocated grid. Station coordinates identify the relevant grid cell; the same spatial indices are used for AGRI, geometry, and COT. The COT retriever uses a $16\times16$ input patch, while the GHI network reads the central $3\times3$ neighbourhood. These are distinct spatial supports: the retrieved local COT can carry information from the wider retrieval patch.

CPP supplies COT and cloud-phase labels for supervised retrieval and reference-based diagnostics. It is a retrieval product with its own uncertainty. Missing or invalid CPP values are excluded from reference-dependent calculations, whereas the main paired GHI cohort is defined without requiring CPP availability at every target.

\subsection{Interval targets and solar geometry}
Let $\tau$ denote the end of a 15-minute averaging interval. The Hunan observations are five-minute interval means stamped at their endpoints, so
\begin{equation}
\bar G_{s,\tau}=\frac{1}{3}\sum_{j=0}^{2}G^{(5)}_{s,\tau-5j\,\mathrm{min}},
\qquad
\bar G^{cs}_{s,\tau}\simeq
\frac{1}{3}\sum_{j=0}^{2}G^{cs}_{s,\tau-(2.5+5j)\,\mathrm{min}} .
\label{eq:interval}
\end{equation}
For Hunan, the clear-sky term uses the Ineichen--Perez model implemented in pvlib, sampled at the midpoints of the three five-minute intervals \cite{ineichen,pvlibclear}; Hubei uses its separately documented one-minute records and masks incomplete target windows. The resulting target and reconstruction are
\begin{equation}
k_{s,\tau}=\bar G_{s,\tau}/\bar G^{cs}_{s,\tau},
\qquad
\widehat G_{s,\tau}=\widehat k_{s,\tau}\bar G^{cs}_{s,\tau}.
\label{eq:kt}
\end{equation}
The clear-sky index is not capped at one. Ground irradiance is used only as a training target or evaluation label. The main cohort requires complete nonnegative GHI records, a positive interval clear-sky estimate, and solar zenith below $90^\circ$ at the station's nominal target time. A sensitivity calculation uses $\bar G^{cs}\geq100$ W m$^{-2}$. Hubei target windows with missing one-minute bins are excluded rather than interpolated.

The spatial geometry features are the cosine of solar zenith, $\cos\theta_0$, and the cosine of the relative solar--viewing azimuth, $\cos\phi_r$. The retriever additionally uses a daytime flag. Geometry is attached to the satellite frame, whereas GHI represents the full preceding interval. Thus, synchronous estimation denotes an association with the same nominal satellite target, not a measurement of instantaneous GHI at the scan time.

\subsection{Temporal partitions and sample definitions}
The Hunan partition reserves separate consecutive two-day validation and test blocks within each sufficiently covered month, with a one-day buffer on either side of each block. Remaining dates are used for training; months with fewer than twenty eligible days are training-only. Normalization uses training dates alone. Retrieval training and GHI evaluation follow the same regional date mapping.

Synchronous examples are unique $(s,\tau)$ pairs. The Hunan GHI dataset contains 50,458 training and 4,479 validation examples, with 2,128 validation examples at Sili and 2,351 at Zhujia. Forecast examples additionally carry initialization time $t$: eight observed frames ending at $t$ predict targets $\tau_k=t+k\Delta$, for $k=1,\ldots,16$ and $\Delta=15$ minutes. Their history and target frames must all belong to the same partition. Sample support is reported separately for these tasks because a valid synchronous example need not belong to a complete forecast sequence. All numerical results currently reported below are validation results.
\FloatBarrier
\section{Methodology}
The method maps spectral observations to a cloud-property representation and then to station irradiance (Fig.~\ref{fig:framework}).

\begin{figure}[!htbp]
\centering
\includegraphics[width=\linewidth]{figures/fig2_cot_estimation_forecast_workflow.pdf}
\caption{COT-assisted station irradiance workflow and evaluation tasks. In synchronous estimation, a frozen retriever maps a $16\times16$ AGRI--geometry patch to a COT field; the GHI head fuses the aligned central $3\times3$ AGRI, geometry, and COT neighbourhood. In the forecasting pathway, fixed SimVP first extrapolates the regional AGRI sequence, after which the same retrieval mapping produces forecast-derived COT. Synchronous estimation and forecast-input evaluation are assessed separately; both pathways produce 15-minute interval-mean GHI.}
\label{fig:framework}
\end{figure}
\subsection{COT retrieval and local representation}
For an observed frame $A^{obs}_\tau$, the retriever estimates a nonnegative log-COT field,
\begin{equation}
L^{obs}_\tau=R_\theta(A^{obs}_\tau,g_\tau),\qquad
C^{obs}_\tau=\exp(L^{obs}_\tau)-1 .
\label{eq:retrieval}
\end{equation}
The input to $R_\theta$ has shape $16\times16\times16$ in channel--row--column order: thirteen AGRI channels and three geometric/daytime channels. The model is a two-downsampling-level U-Net with feature widths 32, 64, and 128 at spatial sizes $16\times16$, $8\times8$, and $4\times4$. Each block contains two $3\times3$ bias-free convolutions, each followed by batch normalization and SiLU. Two max-pooling operations downsample the encoder features. Kernel-2, stride-2 transposed convolutions restore the resolution; the matching encoder features are concatenated at each decoder stage, whose output widths are 64 and 32. A final $1\times1$ convolution maps 32 features to one channel, and softplus enforces $L\geq0$.

The Hunan retriever is initialized randomly and trained against $\log(1+C^{CPP})$. Its loss averages the Huber penalty with transition $\delta=0.1$ over valid pixels,
\begin{equation}
\mathcal L_R=
\frac{\sum_{i,p}m_{i,p}\,
\rho_{0.1}\!\left(L_{i,p}-\log(1+C^{CPP}_{i,p})\right)}
{\sum_{i,p}m_{i,p}},
\label{eq:retrieval-loss}
\end{equation}
where $p$ indexes a pixel and the denominator is required to be positive. The mask retains finite CPP COT in $[0,100]$, accepted CPP phase classes 0--2, finite geometry, valid measurements in all thirteen channels, and daytime pixels with solar zenith below $80^\circ$. Valid zero-COT references participate in training. Retrieval-channel normalization is fitted to valid training support, with the binary daytime channel left unchanged. The Hunan retriever uses seed 42, batch size 1,284, AdamW with learning rate $10^{-3}$ and weight decay $10^{-4}$, gradient clipping at 5, and mixed-precision training. Training runs for at most 50 epochs with patience 8 on validation masked Huber loss. The selected retriever is frozen before GHI training; the subsequent three-seed experiments repeat GHI-head training, not retriever training.

Let $\mathcal P_s$ select the station's central $3\times3$ field. The GHI inputs are standardized using fixed training statistics:
\begin{equation}
X_{s,\tau,c}=\frac{\mathcal P_s A^{obs}_{\tau,c}-\mu_c}{\sigma_c},
\qquad
Z_{s,\tau}=\frac{\mathcal P_s L^{obs}_\tau-\mu_L}{\sigma_L}.
\label{eq:standardization}
\end{equation}
The COT statistics are computed over training examples and local pixels, rather than separately for each image. This retains variation in COT magnitude between samples. Invalid AGRI entries are assigned zero after standardization. The two GHI geometry channels retain their physical cosine values. The local COT input is a nine-pixel field; it is not replaced by a spatial mean or a percentile.

\subsection{Station GHI network}
Each synchronous GHI example contains one target-time frame. The common imagery branch concatenates thirteen AGRI channels with two geometry channels and applies two $1\times1$ convolutions:
\begin{equation}
F_A=\sigma\!\left(W_{A,2}*
\sigma\!\left(W_{A,1}*[X,g]\right)\right)
\in\mathbb R^{32\times3\times3},
\label{eq:agri-encoder}
\end{equation}
where $\sigma$ is SiLU and $*$ denotes convolution. The channel widths are $15\rightarrow32\rightarrow32$. The AGRI + COT model adds a separate COT branch,
\begin{equation}
F_C=\sigma(W_C*Z)\in\mathbb R^{4\times3\times3},
\qquad
v=\operatorname{vec}([F_A,F_C])\in\mathbb R^{324}.
\label{eq:cot-encoder}
\end{equation}
The vectorization preserves a fixed pixel order. A multilayer perceptron with widths $324\rightarrow64\rightarrow32\rightarrow1$, SiLU in the hidden layers, and softplus at the output produces $\widehat k$. Spatial positions can contribute differently through this readout; the architecture alone does not establish that cloud-motion or texture information explains its measured gain.

\begin{table}[!htbp]
\centering\small
\setlength{\tabcolsep}{4pt}
\caption{Synchronous station irradiance models. All branches preserve the $3\times3$ layout until vectorization. The parameter counts include only trainable GHI-head parameters.}
\label{tab:heads}
\begin{tabularx}{\linewidth}{Xlllr}
\toprule
Variant & Common branch & Auxiliary branch & MLP widths & Parameters\\
\midrule
AGRI baseline & $15\to32\to32$ & None & $288\to64\to32\to1$ & 22,177\\
AGRI + COT & $15\to32\to32$ & COT $1\to4$ & $324\to64\to32\to1$ & 24,489\\
AGRI + auxiliary features & $15\to32\to32$ & $F_A$: $32\to4$ & $324\to64\to32\to1$ & 24,613\\
\bottomrule
\end{tabularx}
\end{table}

The imagery-only AGRI baseline omits the auxiliary branch. AGRI + auxiliary features uses a learned $32\rightarrow4$ projection of $F_A$ instead of COT, giving a head parameter count within 0.51\% of AGRI + COT. This arm controls for the added capacity of the GHI head.

\subsection{GHI objective and matched training}
Although the head returns a clear-sky index, optimization uses reconstructed physical GHI:
\begin{equation}
\mathcal L_H=\frac{1}{N}\sum_{i=1}^{N}
\left(\frac{\widehat k_i\bar G^{cs}_i-\bar G_i}
{1000\ \mathrm{W\,m^{-2}}}\right)^2 .
\label{eq:ghi-loss}
\end{equation}
All included examples have unit weight. Separate station heads are trained for seeds 42, 43, and 44 with 30 complete epochs, AdamW learning rate $10^{-3}$ and weight decay $10^{-4}$, cosine annealing to $10^{-5}$, and gradient-norm clipping at 5. Hunan synchronous heads use batch size 2,000; Hubei synchronous heads use 512; Hubei forecast-input heads use 1,536. The checkpoint with the lowest station validation GHI RMSE is retained.

For each station and seed, matched heads share the imagery-branch initialization and corresponding readout weights. Additional readout columns are initially zero, so the auxiliary branch begins as a correction to the same base function. The same epoch-wise sample permutations and optimization budget are used for every head. The COT retriever remains frozen throughout.

\subsection{Fixed extrapolation and downstream error propagation}
The forecasting pathway applies a region-specific, fixed SimVP-family model \cite{simvp,simvpv2} to eight historical frames and deterministic geometry for the history and sixteen target times. The Hunan experiments use the Hunan predictor; Hubei uses its separately trained Hubei predictor and frozen COT retriever:
\begin{equation}
\widehat A_{t+1:t+16}=S(A_{t-7:t},g_{t-7:t+16}),
\qquad
\widehat L_{t+k}=R_\theta(\widehat A_{t+k},g_{t+k}).
\label{eq:forecast}
\end{equation}
Here the indices advance in 15-minute steps. The regional input is $256\times256$. The predictor uses two shared eight-frame prediction steps: the first consumes observed AGRI and historical geometry, and the second consumes the first eight predictions and their geometry. A shared $1\times1$ projection from three geometry channels to thirteen AGRI channels adds target-time conditioning to each predicted segment. Spatial and temporal hidden widths are 32 and 256, with four spatial stages and eight temporal blocks. The Hunan predictor was trained from random initialization for twelve epochs using valid-pixel normalized AGRI MSE; checkpoint selection used season-equal validation pixel RMSE. Hubei forecast-head training used five-fold date-grouped out-of-fold predicted-AGRI/COT inputs; Hubei validation used the frozen full-training imagery and COT predictions.

Predictive geometry is known from time and location and does not require future cloud observations. Predicted AGRI is inverse-standardized to physical channel values and then transformed with the retriever's own normalization. The same local crop and GHI normalization are applied.

The forecasting comparison uses predicted AGRI alone, predicted AGRI with COT, and predicted AGRI with an additional image-feature branch. All models are trained on forecast inputs and evaluated at identical stations, initialization times, and target times. The image-feature branch projects 32 channels to four, yielding 24,613 trainable parameters compared with 24,489 for the COT model. A separate sensitivity analysis holds the synchronous GHI model fixed while replacing AGRI and COT inputs individually with their forecast counterparts. This distinguishes the response to each input from the performance of models trained for forecasting.

\FloatBarrier
\section{Experimental Design}
\subsection{Paired GHI evaluation}
The primary comparison is AGRI + COT against the AGRI baseline on identical station--target samples, observations, and daytime support. For model $m$, seed $r$, station $s$, and common cohort $\mathcal I_s$, the reported regional score is
\begin{equation}
E_m=\frac{1}{N_rN_s}\sum_{r=1}^{N_r}\sum_{s=1}^{N_s}
\sqrt{\frac{1}{|\mathcal I_s|}\sum_{i\in\mathcal I_s}
(\widehat G_{m,r,i}-\bar G_i)^2},
\qquad
J=100\,\frac{E_{\mathrm{base}}-E_{\mathrm{COT}}}{E_{\mathrm{base}}}.
\label{eq:endpoint}
\end{equation}
Here the subscripts base and COT denote the AGRI baseline and AGRI + COT, respectively; $N_r=3$ and $N_s$ is the number of evaluated stations in the region. $J$ is the relative reduction in RMSE, not a percentage-point change in an accuracy measure. RMSE is computed within each seed and station before averaging; predictions are not ensembled before scoring. MAE and signed bias supplement RMSE.

Unless noted otherwise, uncertainty is estimated with 2,000 paired bootstrap resamples of BJT target-date blocks, using the same date weights for every method, station, and seed. The Hunan primary analysis uses 52 distinct dates. The Hubei synchronous primary and consecutive two-day-block sensitivity intervals use 10,000 resamples; the retrieval-feature comparison on common samples uses 2,000 consecutive two-day-block resamples. Forecast-input GHI intervals use 10,000 resamples for both one-day and consecutive two-day blocks. Checkpoints were selected on validation data, so the intervals describe validation cohorts conditional on selection.

For a secondary solar-resource assessment in Hunan, we report GHI errors at station-centre solar altitude $\geq 45^\circ$ on common validation dates. Daily irradiation is calculated as $E_d=0.25\sum_{t\in d}G_t/1000$ in kWh m$^{-2}$, where $G_t$ is the 15-minute interval-mean GHI in W m$^{-2}$. Only dates with every expected 15-minute daylight target at both stations are retained; missing intervals are not interpolated. Date-block bootstrap intervals preserve paired methods and both stations. This radiometric assessment is not an estimate of photovoltaic generation.

As a stronger secondary control, a signed-residual head was trained with matched AGRI-only and AGRI + COT arms. Both arms had 36,773 trainable parameters, zero-initialized corrections, three seeds, 30 epochs, and batch size 2,000. Training used five-fold blocked out-of-fold baseline GHI predictions; validation used the full-training baseline prediction. This comparison is reported separately because the baseline prediction distributions differ between training and validation.

\subsection{Optical-thickness-dependent gains}
The thickness analysis uses the station-centre pixel and the pre-existing COT bins $[0,0.1)$, $[0.1,5)$, $[5,30)$, and $[30,100]$ for valid CPP references. These are operational near-zero, thin, moderate, and thick categories, rather than universal cloud-type definitions. CPP near-zero COT is not equated with an independent clear-sky classification. Each head is scored on precisely the same rows within a bin, using Eq.~(\ref{eq:endpoint}), and counts are reported by station and target date. No whole-patch mean is used to assign thickness.

The reference-based analysis includes only valid centre-CPP matches and is accompanied by an all-cohort analysis grouped by the frozen retriever's centre COT. The latter describes model-retrieved thickness, which can itself be erroneous. Bins are not chosen from observed GHI gains. Subgroup intervals are descriptive and are not adjusted for multiple comparisons; differences between subgroup point estimates do not by themselves establish a statistically resolved trend across thickness.

\subsection{Forecasting and physical cases}
The forecast application is separated into two questions: how well the extrapolated AGRI fields retain COT information, and whether a GHI head trained on forecast inputs can use that information. A fixed-head replacement diagnostic measures sensitivity to forecast AGRI and forecast-derived COT, but is not a deployable forecast score. Forecast-COT scores are computed only on time- and location-matched CPP pixels; CPP is a sparse retrieval reference rather than independent physical truth. The matched forecast-input GHI question remains separate.

Separately trained forecast-input heads were compared on identical station--initialization--target triples using predicted AGRI with or without retrieved COT. This is a validation-only application experiment; the independent test remains unopened. For the synchronous event analysis, four validation sequences were selected without access to GHI values or model errors: one maximum-reference-COT-change event and one reproducibly sampled event per station. The four selected summer cases satisfy the complete 17-frame CPP-reference requirement in the archived product. The event panels distinguish CPP reference COT from the frozen retriever's COT supplied to the GHI head; their GHI curves use observed AGRI and are same-time diagnostics, not forecasts. Neither future-observation replacement nor reference-mask intervention is treated as a deployable forecasting input.
\FloatBarrier
\section{Results}
\subsection{Retrieval quality of the COT input}
CPP is used as a retrieval reference for COT diagnostics, not as independent physical truth. Forecast COT is evaluated separately below on the available CPP-referenced support.

On the common valid pixels, the Hunan retriever had a pooled-pixel patch RMSE of 9.074 COT, MAE of 4.646, and bias of $+0.279$ across 356,169 pixels from 1,398 station--time patches on 46 BJT dates. Its station-equal RMSE was 9.053. The Hubei shared-patch retriever had an RMSE of 13.357, MAE of 9.465, and bias of $-1.303$ across 66,304 pixels from 259 frames on 16 BJT dates. These regional scores use different patches and CPP reference samples and are not a controlled ranking of regional model quality. In both regions, COT was underestimated in the CPP-reference $\geq30$ stratum: the mean bias was $-8.72$ (95\% date-block interval $[-11.53,-5.59]$; 50,924 pixels on 34 dates) in Hunan and $-14.11$ ($[-19.49,-9.02]$; 15,869 pixels on 12 dates) in Hubei. These thickness-conditioned results are descriptive, with sparse date support in Hubei.

\begin{figure}[!htbp]\centering
\includegraphics[width=\linewidth]{figures/fig3_retrieval_quality.pdf}
\caption{Retrieved COT compared with the FY-4B CPP retrieval reference on common valid pixels. The top panels show binned pixel density and the 1:1 line; annotations give pooled-pixel RMSE and bias. The lower panels show retrieval bias within CPP-reference COT bins, with 95\% intervals from 2,000 resamples of BJT target dates. Tick labels give valid pixel count ($N$) and target-date count ($D$). Hunan uses two station-centred 16$\times$16 patches; Hubei uses one shared 16$\times$16 patch covering three stations. CPP is a retrieval reference rather than independent physical truth, and the regional panels are not directly comparable because their spatial and supervision supports differ.}\label{fig:retrieval}\end{figure}

\FloatBarrier
\subsection{Overall irradiance estimation performance}
On 4,479 common validation examples from Sili and Zhujia, the AGRI baseline and AGRI + COT obtained three-seed, station-equal RMSEs of 123.980 and 120.240 W m$^{-2}$. Adding COT reduced RMSE by 3.740 W m$^{-2}$, or 3.02\%, with a paired 95\% date-block interval of $[-6.180,-1.597]$ W m$^{-2}$ for the difference. Both stations improved under each of the three random seeds. Thus, the added representation improves same-time GHI estimation in the Hunan validation comparison.

In the separately trained Hubei validation cohort, three-seed, station-equal RMSE was 74.333 W m$^{-2}$ for AGRI, 70.993 W m$^{-2}$ for AGRI + COT, and 73.839 W m$^{-2}$ for the additional-capacity control. The AGRI + COT minus AGRI point difference was $-3.340$ W m$^{-2}$ (10,000 paired date-block resamples; 95\% interval $[-5.800,-1.118]$); the difference from the capacity control was $-2.846$ W m$^{-2}$ (interval $[-5.147,-0.717]$). AGRI + COT had lower validation RMSE than AGRI and the capacity control in eight of nine station-by-seed point comparisons. A consecutive two-day-block sensitivity gave intervals of $[-6.139,-0.729]$ and $[-5.504,-0.349]$ W m$^{-2}$ for these two comparisons, respectively. 

Figure~\ref{fig:gain} shows three-seed station means and paired date-block intervals for both regions. Values are not pooled across regions because ground-label time sampling differs.

\begin{figure}[!htbp]\centering
\includegraphics[width=\linewidth]{figures/fig4_cot_incremental_value.pdf}
\caption{Incremental value of COT for synchronous GHI estimation. (a) Three-seed mean station RMSE for the AGRI baseline and AGRI + COT; circles and diamonds denote the AGRI baseline and AGRI + COT, respectively, and are connected within station. (b) Within-region, station-equal RMSE reduction (AGRI baseline minus AGRI + COT), with paired validation date-block 95\% intervals for both regions. Sample counts are shown in panel (a) by station and panel (b) by region. Regional scores are reported separately because the ground-label sampling procedures differ; they are not pooled.}\label{fig:gain}\end{figure}

The Hunan increment was also present at station-centre solar altitudes of $45^\circ$ or higher: three-seed, station-equal RMSE was 154.111 W m$^{-2}$ for AGRI and 148.645 W m$^{-2}$ for AGRI + COT (1,527 station--target examples over 35 dates; paired difference $-5.465$ W m$^{-2}$, 95\% date-block interval $[-8.990,-2.152]$). On the 30 dates with complete daylight coverage at both Hunan stations, station- and seed-equal daily irradiation MAE decreased from 0.5491 to 0.5032 kWh m$^{-2}$ (difference $-0.0459$, 95\% interval $[-0.0861,-0.0067]$). The daily bias shifted from $-0.0164$ to $-0.0561$ kWh m$^{-2}$, indicating a more negative signed error despite the lower absolute error. The capacity control remained close to AGRI alone (153.854 W m$^{-2}$ high-sun RMSE; 0.5479 kWh m$^{-2}$ daily irradiation MAE). These Hunan validation diagnostics do not imply a photovoltaic generation gain and are not pooled with Hubei.
\subsection{Gain by reference optical thickness}
Valid station-centre CPP matches covered 1,284 of the 4,479 GHI validation examples and 46 target dates. Figure~\ref{fig:thickness} shows the Hunan thickness-stratified comparison; Table~\ref{tab:thickness} gives the paired values. The bins use the station-centre CPP COT and the same rows for both GHI models. The thick-reference group contained 173 examples across 25 dates. Its RMSE decreased from 74.929 to 58.105 W m$^{-2}$, a reduction of 16.824 W m$^{-2}$ or 22.45\%; all three seeds improved. The lower-thickness groups had smaller point-estimate reductions, approximately 1.5--1.7\%, with intervals spanning zero.



The thick-reference improvement varied by station: Sili RMSE decreased from 69.870 to 64.661 W m$^{-2}$, while Zhujia decreased from 79.987 to 51.549 W m$^{-2}$. Positive GHI bias was reduced from 14.460 to 1.648 W m$^{-2}$ at Sili and from 48.356 to 19.432 W m$^{-2}$ at Zhujia. These results are consistent with a reduction of irradiance overestimation within this reference-thick subgroup. They do not imply a uniform correction for all clouds or establish a causal thickness-response law.

Grouping all 4,479 examples by the retriever's centre COT gave relative reductions of 0.33\%, 1.00\%, 6.09\%, and 24.32\% for near-zero, thin, moderate, and thick retrieved COT. This supplementary grouping uses model outputs rather than reference thickness and is interpreted accordingly. Subgroup gains cannot be averaged directly to reconstruct the overall RMSE gain.


\begin{figure}[!htbp]\centering
\includegraphics[width=\linewidth]{figures/fig5_hunan_cpp_cot_strata.pdf}
\caption{Same-time GHI error by station-centre CPP reference COT in Hunan. (a) Station- and seed-equal RMSE for AGRI baseline and AGRI + COT; labels give the common sample count and number of BJT target dates. (b) Paired date-block difference, $\Delta=$ AGRI + COT minus AGRI baseline, with unadjusted pointwise 95\% percentile intervals from 2,000 resamples; negative values favour AGRI + COT. The pre-existing bins use the station-centre pixel, not a patch mean. This figure reports the Hunan cohort only.}\label{fig:thickness}\end{figure}

\subsection{Cloud-event analysis}
The event comparison links observed AGRI false-colour composites, CPP reference COT, frozen-R COT, and interval-mean GHI for four summer validation sequences with sixteen target times. Across these selected cases, the three-seed mean of the within-sequence RMSE was 213.851 versus 218.408 W m$^{-2}$ (AGRI baseline versus AGRI + COT) for the Sili maximum-change case, 151.268 versus 150.397 for the Sili reproducibly sampled case, 129.933 versus 138.871 for the Zhujia maximum-change case, and 126.700 versus 136.617 for the Zhujia reproducibly sampled case. Thus, COT slightly reduced error in one selected sequence and increased it in three; these descriptive case scores have no event-level uncertainty interval and are not independent samples.

In the Sili maximum-change case (9 July 2025, 13:15 BJT), CPP and R station-pixel COT both rose and later declined, while the GHI heads changed more smoothly than the observed 15-minute means during a rapid irradiance drop. This juxtaposition is consistent with unresolved timing or amplitude errors in the station mapping, but a single event cannot identify their cause. The selected cases therefore illustrate event-level heterogeneity rather than explain the aggregate COT gain. They cover summer only because the archived CPP product lacks complete 15-minute reference sequences in the other seasons.

\begin{figure}[!htbp]\centering
\includegraphics[width=0.98\linewidth]{figures/fig6_sili_high_change_event.png}
\caption{Error-blind, high-COT-change synchronous validation case at Sili (9 July 2025, 13:15 BJT initialization). The four image columns show observed AGRI false-colour composites and, separately, CPP-reference and frozen-retriever COT on the same $16\times16$ patch at +15, +60, +120, and +240 min; the cross marks the station pixel. The lower-left panel compares observed 15-min interval-mean GHI with AGRI-baseline and AGRI + COT predictions using same-time observed AGRI, while the lower-right panel shows station-pixel COT trajectories. Across the 16 validation targets, the seed-equal mean within-case RMSE was 213.851 W m$^{-2}$ for AGRI and 218.408 W m$^{-2}$ for AGRI + COT. Shading denotes the range across seeds, not a confidence interval. This is a synchronous diagnostic with observed imagery, not a forecast result.}\label{fig:events}\end{figure}

\FloatBarrier
\subsection{Role of the COT representation}
The additional-capacity control obtained an RMSE of 123.822 W m$^{-2}$, close to the AGRI baseline. AGRI + COT improved on this control by 3.582 W m$^{-2}$ (95\% interval $[-5.977,-1.479]$). The gain therefore exceeds that obtained by adding a comparably sized imagery-derived branch in this experiment.

We also compared an additional arm that replaces the explicit COT input with a similarly sized local features from the frozen retriever's 16-channel decoder feature map. On the 486 target times shared across all three stations (1,458 station-target rows) over 21 dates, RMSE was 72.867 W m$^{-2}$ for AGRI + COT and 72.785 W m$^{-2}$ for AGRI + retrieval features (COT minus retrieval features $=+0.082$ W m$^{-2}$; 2,000 consecutive two-day-block resamples; 95\% interval $[-1.490,+1.713]$). The Hubei AGRI + COT and retrieval-feature heads had 24,489 and 24,549 trainable parameters, respectively. This interval spans zero, so the validation evidence does not distinguish explicit COT from the retriever-feature control.

A separate, parameter-matched signed-residual experiment provides a stronger AGRI comparison. The AGRI residual model achieved 111.277 W m$^{-2}$, while its AGRI + COT counterpart achieved 113.419 W m$^{-2}$. The COT-minus-AGRI difference was $+2.142$ W m$^{-2}$ (95\% paired date-block interval $[-0.688,+4.745]$), which does not resolve an incremental gain or loss. Both arms improved over the frozen baseline, but the added COT did not improve the point estimate. This is a validation-only same-time retrieval comparison; training uses fold-specific out-of-fold baseline predictions whereas validation uses the full-data baseline prediction, a documented distribution shift that limits interpretation.
\begin{table}[!htbp]\centering\small
\caption{Capacity control on the Hunan synchronous validation cohort. RMSE averages stations and random seeds equally. The difference is the control minus AGRI + COT, in W m$^{-2}$, with a paired shared-date-block 95\% interval.}\label{tab:controls}
\begin{tabularx}{\linewidth}{Xrrl}\toprule
Variant & RMSE & Difference & 95\% interval\\\midrule
AGRI + COT & 120.240 & 0 & Reference\\
AGRI + auxiliary features & 123.822 & 3.582 & $[1.479,5.977]$\\
\bottomrule
\end{tabularx}\end{table}

\subsection{Application to short-term forecasting}
The COT advantage in synchronous estimation did not carry over to the matched Hubei forecasting comparison. Across three stations, three random initializations, and 22 target dates, GHI RMSE was 77.150 W m$^{-2}$ for forecast AGRI alone, 78.155 for forecast AGRI with COT, and 78.980 for the additional image-feature control. The COT-minus-imagery difference was $+1.005$ W m$^{-2}$ (95\% paired date-block interval $[-0.722,+3.780]$). Differences for the image-feature control relative to imagery alone and relative to COT were $+1.830$ ($[-0.517,+4.778]$) and $+0.825$ ($[-2.657,+4.402]$) W m$^{-2}$, respectively. None of these comparisons resolved an improvement; the intervals also do not establish equivalence. A two-day-block analysis gave a similar interval for the COT comparison, $[-0.681,+3.813]$ W m$^{-2}$.

Forecast COT nevertheless contained information about cloud evolution. In Hunan, SimVP followed by COT retrieval achieved a station-and-lead-averaged RMSE of 10.664 on the available local $3\times3$ CPP reference pixels, compared with 20.487 for last-frame persistence. Thick-cloud CSI was 0.423 and 0.306, respectively (Fig.~\ref{fig:forecast}). Equal-budget MSE, temporal-change, and spatial-gradient fine-tuning gave RMSEs of 10.750, 10.756, and 10.759. Direct forecasting of historical COT also failed to improve these scores; the exploratory comparisons and cloud-event diagnostics are reported in Appendix~\ref{sec:forecast-supp}.

Input-replacement analysis helps explain why COT forecast skill need not translate into GHI improvement. For a Hunan model trained on observed inputs, RMSE over 2,817 sequences and 52 dates was 118.14 W m$^{-2}$ with observed AGRI and its retrieved COT, 125.42 when only COT was forecast-derived, 126.46 when only AGRI was forecast, and 139.39 when both inputs were forecast-derived. Replacing COT increased RMSE by 7.28 W m$^{-2}$ (95\% interval $[4.52,10.54]$) with observed AGRI and by 12.93 ($[9.73,16.58]$) with forecast AGRI. These hindsight comparisons reveal sensitivity to input error and distribution shift; they do not isolate the cause of the matched forecasting result in Hubei.

\begin{figure}[!htbp]\centering
\includegraphics[width=0.94\linewidth]{figures/fig7_forecast_cot_supervision.pdf}
\caption{COT forecast error and thick-cloud detection. (a,b) RMSE in the local $3\times3$ neighbourhood at Sili and Zhujia; (c,d) thick-cloud CSI ($\mathrm{COT}\geq30$). Curves use common CPP reference pixels to compare original SimVP, three equal-budget fine-tuning variants, and last-frame persistence. Overall RMSE averages stations and leads equally; CSI pools confusion counts. COT is dimensionless.}\label{fig:forecast}\end{figure}

\FloatBarrier
\section{Discussion}
Retrieved COT improved the compact GHI estimator in both regions, but the comparison with richer representations shows that its value depends on the predictor. A similarly sized image-derived branch did not reproduce the full gain, whereas features from the COT retriever performed comparably to explicit COT in Hubei. The stronger AGRI residual model also showed no additional benefit from COT. Together, these results suggest that cloud-product supervision can provide a useful representation for a compact estimator, without establishing that the scalar optical-thickness output is superior to the information retained in learned features.

The largest Hunan improvement occurred in the CPP thick-cloud subset, where adding COT reduced positive GHI bias. This is consistent with better representation of cloud attenuation, although cloud thickness alone does not determine station irradiance. Cloud position relative to the solar beam, diffuse radiation, and changes within the 15-minute averaging interval can also affect the relationship. The event examples show that even a pronounced COT change does not guarantee a better GHI trajectory. Moreover, the retriever sees a $16\times16$ patch whereas the GHI imagery branch sees $3\times3$; the comparison therefore does not separate cloud-property supervision from wider spatial context.

For solar-resource assessment, the lower Hunan daily irradiation MAE is encouraging, but the more negative daily bias shows that an improvement in absolute error can coexist with systematic underestimation. Forecasting introduces a further limitation: predicted COT outperformed persistence, yet adding it did not improve the matched GHI forecasts. This difference is consistent with sensitivity to errors in cloud evolution and to a changed input distribution. The present comparisons cannot attribute it uniquely to the extrapolation model, the retriever, or their interaction with the GHI estimator.

The study is limited to five stations in two separately trained regions. CPP is a retrieval reference with sparse temporal coverage, and the cloud-event examples cover summer only. In addition, model selection and the reported comparisons use validation data; independent testing is needed to assess generalization. These limitations delimit the result: COT is useful within the tested compact estimation framework, but its incremental value in stronger predictors and forecasting remains unresolved.

\section{Conclusion}
Adding retrieved COT to a compact satellite GHI estimator reduced validation RMSE by 3.02\% in Hunan and 4.49\% in Hubei. In Hunan, the largest reduction occurred in the CPP thick-cloud subset, and daily irradiation MAE decreased over complete observation days. However, COT did not improve the stronger AGRI residual model or outperform matched retrieval features. Its benefit therefore depends on how the satellite information is represented. The corresponding forecasting experiment showed no resolved GHI improvement despite COT skill relative to persistence. The results support COT-assisted synchronous estimation in this setting and identify the retention of that benefit under forecast-input errors as an open problem.


\FloatBarrier
\clearpage
\appendix
\section{Supplementary thickness-stratified values}
\renewcommand{\thetable}{S\arabic{table}}
\setcounter{table}{0}
\begin{table}[!ht]
\centering\small
\caption{Synchronous GHI error stratified by valid CPP COT at the station-centre pixel. Counts list Sili/Zhujia and are fixed across seeds. RMSE and its difference are in W m$^{-2}$; all intervals are descriptive paired date-block 95\% intervals for $\Delta=$ AGRI + COT minus the AGRI baseline.}
\label{tab:thickness}
\setlength{\tabcolsep}{4pt}
\begin{tabularx}{\linewidth}{lrrXXrl}
\toprule
CPP COT & Counts & Dates & AGRI baseline & AGRI + COT & $J$ (\%) & $\Delta$ interval\\
\midrule
$[0,0.1)$ & 143/321 & 31 & 159.165 & 156.812 & 1.48 & $[-6.238,1.157]$\\
$[0.1,5)$ & 71/110 & 29 & 165.869 & 163.045 & 1.70 & $[-6.814,1.268]$\\
$[5,30)$ & 316/150 & 36 & 122.509 & 120.430 & 1.70 & $[-7.179,3.030]$\\
$[30,100]$ & 80/93 & 25 & 74.929 & 58.105 & 22.45 & $[-21.862,-12.359]$\\
\bottomrule
\end{tabularx}
\end{table}
\FloatBarrier
\section{Additional cloud-forecast comparisons}\label{sec:forecast-supp}
A separate single-seed screen asked whether forecasting retrieved COT directly from its eight-frame history could improve on the fixed SimVP-forecast-AGRI-to-retriever chain. Both direct models predicted 16 future COT maps over the full $16\times16$ patch. One used future AGRI-derived retriever COT as its training target; the other substituted CPP values where available and used retriever COT elsewhere. They were evaluated on the same 3,140 validation sequences and CPP-common support (46 target dates; 5,224,460 patch pixels). The validation-selected CPP/R hybrid scored 12.407 COT RMSE, versus 12.512 for the R-target model and 10.953 for the fixed SimVP-to-retriever baseline. On the station-centred local $3\times3$ support, the corresponding values were 12.162, 12.275, and 10.664. For the hybrid model, station--lead-macro RMSE stratified by CPP reference thickness was 8.988 for COT below 1, 5.140 for 1--5, 9.594 for 5--30, and 25.668 for COT at least 30 (2,119,268, 575,625, 1,831,804, and 697,763 patch pixels, respectively). The direct models therefore did not improve the point scores, with the largest remaining error in the thick-reference stratum. This is a one-seed validation screen with model selection on the same validation cohort; the point differences are descriptive, not inferential. CPP remains a retrieval reference rather than independent physical truth. The screen evaluates COT only and does not test a GHI head trained on forecast inputs.

To test whether a larger gradient on absolute COT errors could address the thick-reference deficit, we added a fixed physical-amplitude MSE term to the hybrid model's log-space objective, leaving the training data, architecture, seed, and budget unchanged. On the same validation support, the selected model scored 12.473 RMSE on the full patch and 12.209 on the local $3\times3$, slightly worse than the base hybrid (12.407 and 12.162) and still above the fixed SimVP-to-retriever reference (10.953 and 10.664). The COT $\geq30$ stratum improved by only 0.133 and 0.358 RMSE, respectively, relative to the hybrid base, while the overall scores did not improve. This one-seed screen therefore does not support absolute-amplitude weighting as a remedy; it does not distinguish thickness errors from cloud-position or timing errors, and it provides no GHI forecast evidence.

We next decomposed every unique validation episode in which station-pixel CPP COT was at least 30 for two or more consecutive valid 15-minute targets. The seven episodes (four at Sili and three at Zhujia) occurred on only three BJT dates; for each episode we scored the longest-lead validation forecast that covered the complete event. On the local $3\times3$ support, mean event RMSE for SimVP-to-retriever was 19.53 at Sili and 13.48 at Zhujia, compared with 6.87 and 7.38 when the same retriever was applied to observed future AGRI as a hindsight diagnostic. Pooled thick-cloud CSI was 0.149 and 0.000 for SimVP-to-retriever, versus 0.814 and 0.617 for the hindsight diagnostic. The direct CPP/R-hybrid COT forecaster had mean event RMSE 24.47 and 26.79, respectively. The event signatures vary: SimVP-to-retriever detected some Sili episodes but missed the largest one and all three Zhujia episodes. The observed-future-AGRI diagnostic also retains nonzero error and is not deployable. These seven episodes, concentrated on three dates, are descriptive rather than independent statistical evidence; they suggest a possible limitation of forecast imagery without identifying its cause.

To separate forecasted-image evolution from retrieval response, we compared adjacent 15-min COT changes from SimVP-forecast AGRI and observed future AGRI, both processed by the same frozen retriever. Across 52 BJT target dates (423,900 local pixel transitions per station), forecast-to-retrieval changes were correlated with the observed-AGRI retrieval response (Pearson $r=0.533$ at Sili and $0.593$ at Zhujia), but the through-origin slopes were 0.334 and 0.392, indicating attenuated change amplitude relative to this proxy. Relative to last-frame persistence, model-minus-persistence change-RMSE differences were $-1.737$ (95\% date-block interval $[-2.396,-1.212]$) and $-2.184$ ($[-3.076,-1.494]$) at the two stations. Relative to the original SimVP trajectory, temporal-change fine-tuning produced a change-RMSE difference of $-0.046$ at Sili (95\% interval $[-0.085,-0.013]$) and $-0.013$ at Zhujia ($[-0.047,+0.018]$). This is a retrieval-response diagnostic, not an evaluation against physical COT truth, because both trajectories are produced by the same retriever. On the much smaller support where CPP was valid at both adjacent leads (seven BJT dates per station), forecast-to-retrieval changes had weak association with CPP changes ($r=0.133$ at Sili and $0.019$ at Zhujia), whereas the observed-AGRI retrieval response had higher correlations ($r=0.650$ and $0.608$). This sparse, retrieval-based CPP support cannot identify whether the discrepancy is caused by forecasted cloud evolution, retrieval error, or reference-product variability; these proxy results are mechanistic evidence only.
\FloatBarrier
\begin{thebibliography}{9}\small
\bibitem{ineichen}
P. Ineichen and R. Perez.
A new airmass independent formulation for the Linke turbidity coefficient.
\textit{Solar Energy}, 73:151--157, 2002.
\bibitem{pvlibclear}
pvlib python contributors.
Ineichen/Perez clear-sky model documentation.
\url{https://pvlib-python.readthedocs.io/en/stable/_modules/pvlib/clearsky.html}.
\bibitem{simvp}
Z. Gao, C. Tan, L. Wu, and S. Z. Li.
SimVP: Simpler Yet Better Video Prediction.
In \textit{Proceedings of CVPR}, pp.~3170--3180, 2022.
\url{https://openaccess.thecvf.com/content/CVPR2022/html/Gao_SimVP_Simpler_Yet_Better_Video_Prediction_CVPR_2022_paper.html}.
\bibitem{simvpv2}
C. Tan, Z. Gao, S. Li, and S. Z. Li.
SimVPv2: Towards Simple yet Powerful Spatiotemporal Predictive Learning.
arXiv:2211.12509, 2022 (revised 2024).
\url{https://arxiv.org/abs/2211.12509}.
\bibitem{chen2023}
S. Chen, C. Li, Y. Xie, and M. Li.
Global and direct solar irradiance estimation using deep learning and selected spectral satellite images.
\textit{Applied Energy}, 352:121979, 2023.
\url{https://doi.org/10.1016/j.apenergy.2023.121979}.
\bibitem{sasaki2024}
J. Sasaki, M. Okada, K. Utsunomiya, and K. Yamaguchi.
A Physics-Informed Machine Learning Approach utilizing Multiband Satellite Data for Solar Irradiance Estimation.
arXiv:2407.04283, 2024.
\url{https://arxiv.org/abs/2407.04283}.
\bibitem{pap2024}
K. Papachristopoulou et al.
Effects of clouds and aerosols on downwelling surface solar irradiance nowcasting and short-term forecasting.
\textit{Atmospheric Measurement Techniques}, 17:1851--1877, 2024.
\url{https://doi.org/10.5194/amt-17-1851-2024}.
\bibitem{peng2026}
W. Peng, T. He, Y. Ma, and Y. Zheng.
A physics-informed machine learning method for short-term solar radiation forecast with satellite-derived cloud optical depth.
\textit{Applied Energy}, 406, 127284, 2026.
\url{https://doi.org/10.1016/j.apenergy.2025.127284}.
\bibitem{drosoglou2026}
T. Drosoglou, A. Natsis, D. Dimitriadis, K. Papachristopoulou, S. Kazadzis, A. Georgakis, K. Kilchhofer, G. Tsoumakas, and V. Amiridis.
From cloud optical thickness to solar radiation: deep spatiotemporal forecasting with geostationary satellite imagery.
\textit{EGUsphere} preprint, 2026, under review for \textit{Atmospheric Measurement Techniques}.
\url{https://doi.org/10.5194/egusphere-2026-5862}.
\bibitem{zhang2026}
F. Zhang et al.
Short-term forecasting of cloud physical properties based on Fourier neural operator method.
\textit{Geophysical Research Letters}, 53, e2025GL119553, 2026.
\url{https://doi.org/10.1029/2025GL119553}.
\end{thebibliography}
\end{document}
